% Regression: a string that declares a crossing keeps its declared species
% hue (LionSR/tenkz#2).
%
% The report: `run' declares species=gauge (hue=source:blue) and a crossing
% against `cutter', and drew black -- the declared hue lost.  The suspected
% mechanism was that a declared species' derived styles are local to the
% declaring group while the crossing surgery's deferred redraw strokes
% outside it, so the style lookup missed and the ink fell back to plain.
% Prelude species were said to survive because their hues resolve through the
% kernel's own resolver instead.
%
% It no longer reproduces.  Nothing pinned it, though, and the event stream
% cannot: the stream records `species=gauge' on the wire whether the ink came
% out blue or black, so only pixels witness this.  Both directions are below,
% and the pair is in the kernel pixel baseline.
%
% k1 is the report's own reproduction, which must ink blue.
% k2 is the same picture with the species dropped, which must ink black --
% without it a fixture that stopped inking blue everywhere would still pass.
\documentclass{standalone}
\usepackage{tenkz}
\begin{document}
\tenkzkernel
\tndeclare{species}{gauge}{hue=source:blue}
\begin{tenkz}[rows={wire,wire,wire}, cols=4, bonds=none]
  \tn[at=(1,2), name=v, skin=box]{V}
  \tn[at=(3,1), skin=none, name=p]{}
  \tn[at=(3,4), name=b, skin=box]{B}
  \tnwire[kind=string, name=cutter]{v.270}{open s}
  \tnwire[kind=string, species=gauge, name=run,
    cross={over at crossing of run and cutter}]{p}{b.180}
\end{tenkz}
\begin{tenkz}[rows={wire,wire,wire}, cols=4, bonds=none]
  \tn[at=(1,2), name=v, skin=box]{V}
  \tn[at=(3,1), skin=none, name=p]{}
  \tn[at=(3,4), name=b, skin=box]{B}
  \tnwire[kind=string, name=cutter]{v.270}{open s}
  \tnwire[kind=string, name=run,
    cross={over at crossing of run and cutter}]{p}{b.180}
\end{tenkz}
\end{document}
